\begingroup
\begin{figure*}[ht]
\centering
\hspace{-.05in}
	\begin{subfigure}[b]{0.32\textwidth}
		\includegraphics[width=\textwidth]{A_PQBS_V_PQBR.png}
		\subcaption[]{random v sequential}
		\label{fig:A_PQBS_V_PQBR}
	\end{subfigure}
\hspace{-.05in}
	\begin{subfigure}[b]{0.32\textwidth}
		\includegraphics[width=\textwidth]{A_PQBR_THROUGH_PUT_ALL.png}
		\subcaption[]{throughput}
		\label{fig:A_PQBR_THROUGH_PUT_ALL}
	\end{subfigure}
\hspace{-.05in}
	\begin{subfigure}[b]{0.32\textwidth}
		\includegraphics[width=\textwidth]{A_PQBR_TIME_PERP_ALL.png}
		\subcaption[]{time per particle}
		\label{fig:A_PQBR_TIME_PERP_ALL}
	\end{subfigure}
\hspace{-.05in}
\caption[TITLE:PQBR Performance Plots]{(a) Comparison of randomly ordered particle occupancy (red squares) and
sequentially ordered particle occupancy (blue circles). Sequential
ordering consistently produces lower execution times throughout the
evaluated particle range. At the largest benchmark configuration,
random ordering increases execution time by approximately
\nc{Locality Penalty} relative to sequential ordering, highlighting
the importance of memory locality despite identical algorithmic
complexity.
(b) Log--log plot of total throughput (red), compute throughput (blue),
and graphics throughput (black) as a function of particle count. The
logarithmic axes permit visualization of performance trends across
nearly three orders of magnitude in particle count. Dashed vertical
lines indicate observed transition regions within the measured
throughput behavior.

A log-log plot of throughput for graphics, compute, and total. The saturated point plateau's after $N=$\numprint{\APQBRTIMEPERPALLstartval{}} particles shown in the plot as a black vertical dashed line. 

(c) time per particle for graphics, compute, total starting at, N=\numprint{\APQBRTIMEPERPALLstartval{}}, for the start of the heavily saturated range. Graphics stabilizes at $\APQBRTIMEPERPALLGraphicsstabilzes{}  \pm \APQBRTIMEPERPALLGraphicsstdval{} $ $\sfrac{ns}{particle}$, compute at $\APQBRTIMEPERPALLComputestabilzes{} \pm \APQBRTIMEPERPALLComputestdval{}$ $\sfrac{ns}{particle}$ for a total stabilization at $\APQBRTIMEPERPALLTotalstabilzes{} \pm \APQBRTIMEPERPALLTotalstdval{}$ $\sfrac{ns}{particle}$. }
\Description[TITLE:PQBR Performance Plots]{(a) A plot of randomly ordered particles performance (red squares) versus
sequentially ordered particles (blue circles). The random performance 
increases disproportionately with sequential performance as the number of 
particles increase.
(b) A log-log plot of throughput for graphics, compute, and total. The saturated point plateau's after $N=$\numprint{\APQBRTIMEPERPALLstartval{}} particles shown in the plot as a black vertical dashed line. (c) time per particle for graphics, compute, total starting at, N=\numprint{\APQBRTIMEPERPALLstartval{}}, for the start of the heavily saturated range. Graphics stabilizes at $\APQBRTIMEPERPALLGraphicsstabilzes{}  \pm \APQBRTIMEPERPALLGraphicsstdval{} $ $\sfrac{ns}{particle}$, compute at $\APQBRTIMEPERPALLComputestabilzes{} \pm \APQBRTIMEPERPALLComputestdval{}$ $\sfrac{ns}{particle}$ for a total stabilization at $\APQBRTIMEPERPALLTotalstabilzes{} \pm \APQBRTIMEPERPALLTotalstdval{}$ $\sfrac{ns}{particle}$. }
\label{fig:A_PQBR_THROUGHPUT_AND_TPP}
\end{figure*}
\endgroup